% =============================================================================
% content.tex — Codara PFE Report
% =============================================================================

\documentclass[12pt,a4paper]{report}

\usepackage[T1]{fontenc}
\usepackage[utf8]{inputenc}
\usepackage{times}
\usepackage[english]{babel}

\usepackage[
  a4paper,
  top=2.5cm, bottom=2.5cm,
  left=3cm,  right=2.5cm,
  headheight=15pt
]{geometry}

\usepackage{setspace}
\onehalfspacing

\usepackage{fancyhdr}
\pagestyle{fancy}
\fancyhf{}
\fancyhead[R]{\nouppercase{\leftmark}}
\fancyfoot[C]{\thepage}
\renewcommand{\headrulewidth}{0.4pt}

\setcounter{secnumdepth}{3}
\setcounter{tocdepth}{3}

\usepackage{titlesec}
\titleformat{\chapter}[display]
  {\normalfont\Large\bfseries\centering}
  {Chapter \thechapter}{10pt}{\LARGE}
\titlespacing*{\chapter}{0pt}{30pt}{20pt}

\usepackage{titletoc}

\usepackage{graphicx}
\usepackage{booktabs}
\usepackage{longtable}
\usepackage{array}
\usepackage{multirow}
\usepackage{listings}
\usepackage{xcolor}

\usepackage{amsmath, amssymb}

\usepackage[hidelinks, breaklinks=true]{hyperref}
\usepackage{bookmark}

\usepackage[numbers, sort&compress]{natbib}
\bibliographystyle{IEEEtran}

\usepackage[printonlyused]{acronym}

\newcommand{\codara}{\textsc{Codara}}

% =============================================================================
\begin{document}

\pagenumbering{roman}

% Title Page
\begin{titlepage}
  \centering
  \vspace*{1cm}
  {\large \textbf{[Institution Name]}}\\[0.4cm]
  {\large [Faculty / Department]}\\[0.5cm]
  \rule{\linewidth}{0.5pt}\\[0.5cm]
  {\Large \textbf{End-of-Study Project Report (PFE)}}\\[0.3cm]
  {\large Engineering Degree in [Specialty]}\\[1.2cm]
  {\LARGE \textbf{\codara{}}}\\[0.5cm]
  {\large \textit{A Multi-Agent SAS-to-Python/PySpark Conversion Accelerator}}\\[0.3cm]
  {\normalsize Using LangGraph Orchestration, RAPTOR Semantic Clustering,}\\
  {\normalsize 3-Tier LLM Fallback Translation, Z3 Formal Verification,}\\
  {\normalsize and SemantiCheck 4-Layer Semantic Verification}\\[2cm]
  \begin{tabular}{ll}
    \textbf{Presented by:} & [Author Name] \\[0.3cm]
    \textbf{Supervisor:}   & [Supervisor Name], [Title] \\[0.3cm]
    \textbf{Host company:} & [Company Name] \\
  \end{tabular}\\[2cm]
  \begin{tabular}{lll}
    \multicolumn{3}{c}{\textbf{Jury}} \\[0.3cm]
    [Name] & [Title] & President \\
    [Name] & [Title] & Examiner \\
    [Name] & [Title] & Supervisor \\
  \end{tabular}\\[2cm]
  {\large Academic Year 2025--2026}
\end{titlepage}

\chapter*{Dedication}
\addcontentsline{toc}{chapter}{Dedication}
\thispagestyle{plain}

\chapter*{Acknowledgements}
\addcontentsline{toc}{chapter}{Acknowledgements}
\thispagestyle{plain}

\chapter*{Abstract}
\addcontentsline{toc}{chapter}{Abstract}
\thispagestyle{plain}

\tableofcontents

\listoffigures
\addcontentsline{toc}{chapter}{List of Figures}

\listoftables
\addcontentsline{toc}{chapter}{List of Tables}

\chapter*{List of Abbreviations}
\addcontentsline{toc}{chapter}{List of Abbreviations}
\thispagestyle{plain}

\begin{acronym}[HyperRAPTOR]
  \acro{API}{Application Programming Interface}
  \acro{AST}{Abstract Syntax Tree}
  \acro{ASGI}{Asynchronous Server Gateway Interface}
  \acro{CDAIS}{Constraint-Driven Adversarial Input Synthesis}
  \acro{CEGAR}{Counterexample-Guided Abstraction Refinement}
  \acro{CI/CD}{Continuous Integration / Continuous Delivery}
  \acro{CRISP-DM}{Cross-Industry Standard Process for Data Mining}
  \acro{DAG}{Directed Acyclic Graph}
  \acro{DSR}{Design Science Research}
  \acro{ECE}{Expected Calibration Error}
  \acro{ETL}{Extract, Transform, Load}
  \acro{FSM}{Finite State Machine}
  \acro{GMM}{Gaussian Mixture Model}
  \acro{IR}{Intermediate Representation}
  \acro{JWT}{JSON Web Token}
  \acro{KB}{Knowledge Base}
  \acro{LLM}{Large Language Model}
  \acro{MRR}{Mean Reciprocal Rank}
  \acro{PFE}{Projet de Fin d'Études (End-of-Study Project)}
  \acro{RAG}{Retrieval-Augmented Generation}
  \acro{RAPTOR}{Recursive Abstractive Processing for Tree-Organized Retrieval}
  \acro{REST}{Representational State Transfer}
  \acro{SAS}{Statistical Analysis System}
  \acro{SCC}{Strongly Connected Component}
  \acro{SCS}{SemantiCheck Score}
  \acro{SDLC}{Software Development Life Cycle}
  \acro{SMT}{Satisfiability Modulo Theories}
  \acro{SSE}{Server-Sent Events}
  \acro{STG}{Semantic Transformation Graph}
\end{acronym}

% =============================================================================
\pagenumbering{arabic}
\setcounter{page}{1}

\chapter*{General Introduction}
\addcontentsline{toc}{chapter}{General Introduction}
\markboth{General Introduction}{General Introduction}
\thispagestyle{fancy}

% =============================================================================
\chapter{Project Context and Foundation}
\label{ch:context}

\section{Introduction}
\label{sec:ch1-intro}

\section{Academic Institution}
\label{sec:institution}

\section{Host Company}
\label{sec:host-company}

  \subsection{Company Profile and Business Context}
  \label{subsec:company-overview}

  \subsection{Mission, Services, and Technological Focus}
  \label{subsec:mission-vision}

  \subsection{Internship Scope}
  \label{subsec:internship-scope}

\section{Project Presentation}
\label{sec:project-presentation}

  \subsection{Initial Situation and Problem Statement}
  \label{subsec:problem}

  \subsection{Study and Critique of Existing Solutions}
  \label{subsec:existing-solutions}

  \subsection{Proposed Solution Overview}
  \label{subsec:proposed-solution}

  \subsection{Expected Deliverables}
  \label{subsec:deliverables}

\section{Methodology Orientation}
\label{sec:methodology}

  \subsection{Overview of Candidate Methodologies}
  \label{subsec:methodologies-overview}

  \subsection{Design Science Research (DSR)}
  \label{subsec:dsr}

  \subsection{Adopted Hybrid Methodology: DSR + Agile}
  \label{subsec:hybrid-methodology}

\section{Conclusion}
\label{sec:ch1-conclusion}


% =============================================================================
\chapter{Theoretical Foundations and Related Work}
\label{ch:foundations}

\section{Introduction}
\label{sec:ch2-intro}

\section{The SAS Language: Constructs and Migration Challenges}
\label{sec:sas-language}

  \subsection{Core Language Constructs}
  \label{subsec:sas-constructs}

  \subsection{Python and PySpark Equivalents}
  \label{subsec:python-equivalents}

  \subsection{Complexity Taxonomy for Migration}
  \label{subsec:complexity-taxonomy}

\section{Automated Code Translation}
\label{sec:code-translation}

  \subsection{Rule-Based Transpilers}
  \label{subsec:rule-based}

  \subsection{LLM-Based Code Translation}
  \label{subsec:llm-translation}

  \subsection{Hybrid Approaches}
  \label{subsec:hybrid-translation}

  \subsection{Limitations of Existing Tools}
  \label{subsec:existing-limitations}

\section{Retrieval-Augmented Generation}
\label{sec:rag-theory}

  \subsection{Foundations of RAG}
  \label{subsec:rag-foundations}

  \subsection{RAPTOR: Recursive Abstractive Processing}
  \label{subsec:raptor-theory}

  \subsection{Graph RAG}
  \label{subsec:graph-rag-theory}

  \subsection{Agentic RAG}
  \label{subsec:agentic-rag-theory}

\section{Multi-Agent Orchestration}
\label{sec:mas-theory}

  \subsection{Multi-Agent System Principles}
  \label{subsec:mas-principles}

  \subsection{LangGraph as Orchestration Framework}
  \label{subsec:langgraph-theory}

  \subsection{Why Multi-Agent for SAS Migration}
  \label{subsec:mas-justification}

\section{Formal Verification for Generated Code}
\label{sec:formal-verification-theory}

  \subsection{Satisfiability Modulo Theories (SMT)}
  \label{subsec:smt-theory}

  \subsection{Code Equivalence Checking}
  \label{subsec:equivalence-checking}

  \subsection{CEGAR: Counterexample-Guided Abstraction Refinement}
  \label{subsec:cegar-theory}

  \subsection{Program Slicing for Fault Localisation}
  \label{subsec:program-slicing-theory}
  % Weiser (1984) backward slicing: given a failing statement,
  % the slice is the minimal subset of statements that affect its value.
  % Applied here to reduce the repair context sent to the LLM after
  % a sandbox execution failure, improving repair precision.

\section{Semantic Equivalence Verification Beyond SMT}
\label{sec:semantic-equiv-theory}

  \subsection{Oracle-Based Behavioural Testing}
  \label{subsec:oracle-testing-theory}

  \subsection{Adversarial Test Synthesis}
  \label{subsec:adversarial-synthesis-theory}

  \subsection{Composite Verification Scoring}
  \label{subsec:composite-scoring-theory}

\section{Machine Learning for Complexity Scoring}
\label{sec:ml-theory}

  \subsection{Logistic Regression and Calibration}
  \label{subsec:lr-calibration}

  \subsection{Feature Engineering for Code Complexity}
  \label{subsec:code-features}

\section{Related Work}
\label{sec:related-work}

  \subsection{Migration Tools, Validation, and Multi-Agent Approaches}
  \label{subsec:related-migration}

  \subsection{Research Gap and Positioning of This Project}
  \label{subsec:research-gap}

\section{Synthesis}
\label{sec:ch2-synthesis}

\section{Conclusion}
\label{sec:ch2-conclusion}


% =============================================================================
\chapter{Project Specifications}
\label{ch:specifications}

\section{Introduction}
\label{sec:ch3-intro}

\section{Project Objectives}
\label{sec:objectives}

  \subsection{Business Objectives}
  \label{subsec:business-objectives}

  \subsection{Technical Objectives}
  \label{subsec:technical-objectives}

  \subsection{Scientific Objectives}
  \label{subsec:scientific-objectives}

\section{Functional Requirements}
\label{sec:functional-requirements}

\section{Non-Functional Requirements}
\label{sec:nonfunctional-requirements}

\section{Use Cases and Scenarios}
\label{sec:use-cases}

  \subsection{Actor Identification}
  \label{subsec:actors}

  \subsection{Use Case Diagram}
  \label{subsec:use-case-diagram}

  \subsection{Detailed Use Cases}
  \label{subsec:detailed-use-cases}

  \subsection{Conversion Scenarios}
  \label{subsec:conversion-scenarios}

\section{High-Level Architecture}
\label{sec:high-level-arch}

\section{Constraints and Assumptions}
\label{sec:constraints}

\section{Conclusion}
\label{sec:ch3-conclusion}


% =============================================================================
\chapter{System Architecture and Design Decisions}
\label{ch:architecture}

\section{Introduction}
\label{sec:ch4-intro}

\section{Technology Selection Through Empirical Evaluation}
\label{sec:tech-selection}

  \subsection{Multi-Agent Orchestration Framework Selection}
  \label{subsec:framework-selection}

  \subsection{LLM Benchmarking and Provider Selection}
  \label{subsec:llm-benchmarking}
  % Chain as of v3.1.0 (April 2026):
  % Primary: Azure GPT-4o (all risk levels)
  % Fallback 1: Ollama nemotron-3-super:cloud
  % Fallback 2: Groq LLaMA-3.3-70b (last resort + cross-verifier)
  % Determined empirically: minimax-m2.7:cloud scored 10/10 on torture_test
  % then nemotron-3-super:cloud adopted; Azure promoted to primary for
  % production reliability; Ollama demoted to warm fallback.

  \subsection{Embedding Model and Vector Store Selection}
  \label{subsec:embedding-selection}

  \subsection{Formal Verification Tool Selection}
  \label{subsec:z3-selection}

\section{The 8-Node LangGraph Pipeline}
\label{sec:pipeline-design}

  \subsection{State Model and Composite Node Design Pattern}
  \label{subsec:state-model}

  \subsection{Input Phase: File Processing and Streaming (Nodes 1--2)}
  \label{subsec:nodes-input}

  \subsection{Partitioning Phase: Boundary Detection and RAPTOR Clustering (Nodes 3--4)}
  \label{subsec:nodes-partition}

  \subsection{Routing, Persistence, Translation, and Merge (Nodes 5--8)}
  \label{subsec:nodes-routing-merge}

\section{Data Architecture}
\label{sec:data-architecture}

  \subsection{SQLite API Database}
  \label{subsec:sqlite-schema}

  \subsection{LanceDB Knowledge Base}
  \label{subsec:lancedb-schema}

  \subsection{DuckDB Audit Store}
  \label{subsec:duckdb-schema}

  \subsection{Redis Checkpoint Store}
  \label{subsec:redis-schema}

  \subsection{Azure Blob and Queue Storage}
  \label{subsec:azure-storage-schema}
  % BlobStorageService: upload/download/list/delete against Azure Blob Storage.
  % Graceful local-disk fallback (backend/uploads/) when not configured.
  % PipelineQueueService: producer + daemon consumer thread; enqueues JSON
  % payloads base64-encoded into Azure Queue (visibility_timeout=300 s);
  % poison-message guard at 5 dequeue attempts.

\section{API and Service Layer Design}
\label{sec:api-design}

  \subsection{Layered Architecture: Routes \texorpdfstring{$\to$}{->} Services \texorpdfstring{$\to$}{->} Core}
  \label{subsec:layered-api}
  % Three service modules extracted in Week 15:
  % pipeline_service.py: run_pipeline_sync() — runs 8 stages, handles
  %   blob download/temp-file cleanup.
  % blob_service.py: BlobStorageService — file I/O abstraction.
  % queue_service.py: PipelineQueueService — async job dispatch.

  \subsection{Key API Endpoints}
  \label{subsec:key-endpoints}

  \subsection{Authentication and Authorisation Design}
  \label{subsec:auth-design}
  % JWT HS256, 24 h expiry. bcrypt password hashing.
  % Rate limiting: 5 attempts / IP / 60 s on /login and /register.
  % Passwords: generated via secrets.token_urlsafe(18) on first boot;
  %   pinned via CODARA_ADMIN_PASSWORD / CODARA_USER_PASSWORD env vars.

\section{Gold Standard Corpus Construction}
\label{sec:corpus-construction}

  \subsection{Motivation and Design Requirements}
  \label{subsec:corpus-design}

  \subsection{Construction Methodology}
  \label{subsec:corpus-methodology}

  \subsection{Statistics and Limitations}
  \label{subsec:corpus-statistics}

\section{Conclusion}
\label{sec:ch4-conclusion}


% =============================================================================
\chapter{Knowledge Integration and Pipeline Implementation}
\label{ch:implementation}

\section{Introduction}
\label{sec:ch5-intro}

\section{Three-Tier RAG Strategy}
\label{sec:rag-impl}

  \subsection{Conceptual Framework, Knowledge Base, and Escalation Rules}
  \label{subsec:rag-framework}

  \subsection{Tier 1: Static RAG}
  \label{subsec:static-rag-impl}

  \subsection{Tier 2: Graph RAG}
  \label{subsec:graph-rag-impl}

  \subsection{Tier 3: Agentic RAG}
  \label{subsec:agentic-rag-impl}

\section{RAPTOR Implementation}
\label{sec:raptor-impl}

  \subsection{Embedding and Clustering Pipeline}
  \label{subsec:embedding-clustering}

  \subsection{Summarisation and Tree Construction}
  \label{subsec:summarizer-impl}

  \subsection{HyperRAPTOR Extension}
  \label{subsec:hyperraptor-impl}

\section{Complexity Scoring and Strategy Routing}
\label{sec:complexity-impl}

  \subsection{Feature Engineering}
  \label{subsec:feature-engineering}

  \subsection{Logistic Regression with Platt Scaling}
  \label{subsec:platt-impl}

  \subsection{StrategyAgent: Risk-to-Paradigm Mapping}
  \label{subsec:strategy-impl}

\section{Translation and Validation}
\label{sec:translation-impl}

  \subsection{Prompt Engineering}
  \label{subsec:prompt-engineering}

  \subsection{LLM Invocation with Structured Output}
  \label{subsec:llm-invocation}

  \subsection{LLM Fallback Chain Implementation}
  \label{subsec:fallback-impl}
  % Chain: Azure GPT-4o (primary) → Ollama nemotron-3-super:cloud (fallback 1)
  % → Groq LLaMA-3.3-70b (fallback 2 + cross-verifier).
  % Circuit breakers: Azure (5 failures → 60 s open),
  %   Groq (3 failures → 120 s open).
  % Rate limiters: Azure (10 concurrent), Groq (3 concurrent).
  % All LLM calls wrapped in asyncio.wait_for(timeout=60 s) for
  % Python 3.10 compatibility (asyncio.timeout() requires 3.11+).

  \subsection{Deterministic Translation Rules}
  \label{subsec:deterministic-rules}
  % deterministic_translator.py: short-circuits LLM for structurally
  % unambiguous SAS constructs. Two new rules added in Week 15+:
  %
  % Rule 1 — PROC FORMAT VALUE (discrete label mappings):
  %   Converts VALUE blocks with string/integer → label entries to a Python
  %   dict. Numeric range formats (18-64, low-high) detected and rejected
  %   to LLM. Output: sexfmt_fmt = {'M': 'Male', 'F': 'Female'} style.
  %
  % Rule 2 — PROC SQL simple SELECT:
  %   Handles single-table SELECT with no subqueries, GROUP BY, or HAVING.
  %   Maps directly to pandas df[cols] or df.query() expressions.
  %
  % Both rules registered in _RULES and tried before LLM dispatch.

  \subsection{Error Analysis and Program Slicing}
  \label{subsec:error-analysis-slicing}
  % error_analyst.py — fault localisation for the repair LLM.
  %
  % Based on Weiser (1984) backward program slicing: given a failing
  % statement, the slice is the minimal subset of statements that
  % causally influence its value.
  %
  % Implementation:
  %   _extract_failing_lineno(): parses "File ..., line X" from tracebacks.
  %   _slice_around_line(): returns ±3 lines with >>> marker.
  %   _slice_for_columns(): AST-walks translated Python for df['col'] /
  %     df.col accesses matching affected column names; expands ±2 lines.
  %
  % Effect: repair prompt contains only the relevant excerpt instead of
  % the full translated script (which may be 200+ lines), improving
  % repair precision and reducing token cost.

  \subsection{ValidationAgent Sandbox}
  \label{subsec:sandbox-impl}
  % multiprocessing.Process isolation: removed builtins (open, __import__,
  % exec, eval, compile, exit, quit, input, breakpoint).
  % Killed via .kill() on timeout (not .join(timeout) which leaks threads).
  % Windows-compatible (no signal.alarm).

  \subsection{Oracle-Based Semantic Validation}
  \label{subsec:semantic-validator-impl}
  % semantic_validator.py: runs translated Python against a reference
  % (oracle) implementation derived from deterministic SAS semantics.
  % Validates: column set, row count, dtype consistency, group aggregates.
  % Complement to syntax sandbox: sandbox catches crashes;
  % semantic validator catches wrong-but-running translations.

  \subsection{Cross-Verification and Reflexion}
  \label{subsec:crossverify-impl}

  \subsection{Lineage Guard}
  \label{subsec:lineage-guard-impl}
  % lineage_guard.py: disk-reload check (original) extended with macro
  % reference check. Flags partitions that reference %macro_name symbols
  % without a corresponding definition in the dependency graph, preventing
  % silent unresolved-reference failures at merge time.

\section{Z3 Formal Verification Layer}
\label{sec:z3-impl}

  \subsection{Z3VerificationAgent Architecture}
  \label{subsec:z3-arch-impl}

  \subsection{SMT Encoding Patterns}
  \label{subsec:z3-patterns}
  % 4 patterns: linear arithmetic, boolean filter, sort/dedup, assignment.
  % Coverage: ~30\% of SAS constructs amenable to SMT encoding.

  \subsection{CEGAR Loop}
  \label{subsec:cegar-impl}

\section{Constraint-Driven Adversarial Input Synthesis (CDAIS)}
\label{sec:cdais-impl}

  \subsection{Motivation and Design Rationale}
  \label{subsec:cdais-motivation}
  % Existing ValidationAgent fuzzes with: empty DataFrame, single-row,
  % all-null column. These catch crashes. They do NOT catch semantic errors:
  % empty DataFrame returns trivial output; single row never fires
  % RETAIN/FIRST./LAST. group boundaries.
  % CDAIS synthesises adversarial inputs specifically tuned to expose the
  % six known SAS→Python migration failure modes.

  \subsection{The Six Adversarial Patterns}
  \label{subsec:cdais-patterns}
  % Six patterns injected into every generated DataFrame:
  % 1. NaN injection (8\% of numeric cells) — SAS missing (.) vs pandas NaN
  %    arithmetic divergence.
  % 2. Currency strings — dtype coercion bugs in LLM-translated code.
  % 3. Multi-group BY layout (3 groups × 10 rows, pre-sorted) — exercises
  %    RETAIN reset and FIRST./LAST. boundaries.
  % 4. Exact duplicate rows (2 extra) — NODUP/NODUPKEY translation bugs.
  % 5. Mixed-case strings (UPPER/lower/Title cycling) — SAS case-sensitive
  %    comparison vs Python.
  % 6. Numeric edge cases (0, negative, 999999.99) — overflow, zero-div,
  %    negative-cumsum.

  \subsection{Z3-Driven Minimum Witness Synthesis}
  \label{subsec:cdais-z3}
  % For structural failure classes (e.g.\ RETAIN\_RESET), Z3 encodes the
  % correct cumsum-per-group and incorrect global-cumsum as symbolic
  % constraints, then minimises Σ values subject to divergence.
  % Result: the smallest concrete DataFrame guaranteed to expose the error
  % class. Differs from z3_agent.py (verification) — here Z3 produces
  % test inputs, not proofs.

  \subsection{Coverage Certificates}
  \label{subsec:cdais-certificates}
  % A translation that passes test T(C) for error class C receives a
  % formal coverage certificate: "free from error class C for any dataset
  % of structural shape S." Heuristic test data cannot provide this
  % guarantee. CDAIS goes from "we tested it" to "we proved it doesn't
  % have this bug."

  \subsection{Column Type Inference}
  \label{subsec:cdais-type-inference}
  % Deterministic regex analysis of SAS source (no LLM):
  % retain / sum / mean / lag() / arithmetic → numeric columns.
  % = "literal" / IN (...) / UPCASE() / INPUT col $ → string columns.
  % BY columns → categorical string groups.

\section{SemantiCheck: 4-Layer Semantic Verification Framework}
\label{sec:semanticheck-impl}

  \subsection{Motivation}
  \label{subsec:semanticheck-motivation}
  % CodeBLEU compares token sequences: two semantically identical
  % translations can score low if styled differently. LLM confidence
  % scores are self-reported and biased. Z3 covers ~30\% of SAS patterns.
  % SemantiCheck composes four orthogonal verification signals into a
  % single \ac{SCS} ∈ [0, 1].

  \subsection{Score Composition}
  \label{subsec:semanticheck-score}
  % SCS = 0.30 × L1 + 0.30 × L2 + 0.20 × L3 + 0.20 × L4
  %
  % Verdicts:
  %   VERIFIED         SCS ≥ 0.85
  %   LIKELY_CORRECT   SCS ≥ 0.65
  %   UNCERTAIN        SCS ≥ 0.40
  %   LIKELY_INCORRECT SCS < 0.40

  \subsection{Layer 1 --- Formal (Z3 SMT)}
  \label{subsec:sc-l1}
  % Weight: 30\%. Reuses Z3VerificationAgent's 4 SMT patterns.
  % Produces: formal\_proof / z3\_unknown / unverifiable.

  \subsection{Layer 2 --- Behavioural (CDAIS Witness Execution)}
  \label{subsec:sc-l2}
  % Weight: 30\%. Runs translated code on the minimum witness DataFrame
  % synthesised by CDAIS for each applicable error class.
  % Pass → 1.0; partial → proportional; fail → 0.0.

  \subsection{Layer 3 --- Contract (Semantic Transformation Graph)}
  \label{subsec:sc-l3}
  % Weight: 20\%. Extracts a language-agnostic \ac{STG} from both
  % the SAS source (regex + heuristics) and the Python translation
  % (AST walk). STG nodes: READ, FILTER, SORT, GROUP, AGGREGATE,
  % COMPUTE, MERGE, DEDUP, TRANSPOSE, WRITE.
  % Score: Jaccard similarity of the two operation-type sets.
  % Syntax-independent: two different implementations of identical logic
  % score identically.

  \subsection{Layer 4 --- Oracle (LLM-as-SAS-Interpreter)}
  \label{subsec:sc-l4}
  % Weight: 20\%. Synthesises a representative 3-row input DataFrame
  % from column names found in the SAS source. Asks the LLM to predict
  % (a) SAS output and (b) Python output on the same input using
  % independent prompts. Scores agreement on:
  %   column set (50\%), row-count direction (30\%), transformation
  %   keywords (20\%).
  % First tool to use LLM-as-SAS-oracle for behavioural testing without
  % a SAS licence.

  \subsection{Orthogonality with Z3 Verification}
  \label{subsec:sc-orthogonality}
  % z3_agent.py covers: linear arithmetic, boolean filter, sort/dedup,
  % assignment (SMT arithmetic patterns).
  % CDAIS + SemantiCheck cover: RETAIN, LAG, SORT stability, NULL
  % arithmetic, JOIN type, GROUP boundaries (dataflow semantic patterns).
  % Together the two systems cover the full semantic surface.

\section{Namespace Safety and Merge-Time Checks}
\label{sec:namespace-impl}

  \subsection{AST Namespace Checker}
  \label{subsec:namespace-checker-impl}
  % namespace_checker.py: AST-level use-before-def and variable shadow
  % checker inserted into MergeAgent before script output.
  % Detects: references to names not yet assigned in the merged script;
  % local assignments that shadow outer-scope variables unexpectedly.

  \subsection{Execution State Tracker}
  \label{subsec:execution-state-impl}
  % execution_state.py: maintains a registry of materialised DataFrame
  % schemas throughout the pipeline. Each node that produces or transforms
  % a DataFrame records its column set, dtypes, and row-count bound.
  % Enables downstream agents (ValidationAgent, namespace checker) to
  % resolve column references without executing code.

\section{Regression Testing Infrastructure}
\label{sec:regression-impl}

  \subsection{Oracle Regression Runner}
  \label{subsec:regression-runner-impl}
  % benchmark/regression_runner.py: CI-safe oracle regression runner.
  % For each gold standard pair: generates adversarial DataFrame via
  % DummyDataGenerator, runs reference Python (gold) and translated
  % Python, compares output DataFrames with tolerance for float drift.
  % Integrated into GitHub Actions pipeline; fails CI on semantic
  % regression (not just syntax errors).

\section{Functional Scenario Testing}
\label{sec:scenario-testing}

  \subsection{Testing Methodology and Scope}
  \label{subsec:test-methodology}

  \subsection{Scenario A: Low-Complexity Static RAG Path}
  \label{subsec:scenario-a}

  \subsection{Scenario B: Cross-File Graph RAG Path}
  \label{subsec:scenario-b}

  \subsection{Scenario C: High-Complexity Agentic RAG + CEGAR}
  \label{subsec:scenario-c}

\section{Conclusion}
\label{sec:ch5-conclusion}


% =============================================================================
\chapter{Deployment and User Interface}
\label{ch:deployment}

\section{Introduction}
\label{sec:ch6-intro}

\section{Backend Deployment}
\label{sec:backend-deployment}

  \subsection{FastAPI Structure, Endpoints, and Background Tasks}
  \label{subsec:fastapi-structure}

  \subsection{Service Layer: Pipeline, Blob, and Queue}
  \label{subsec:service-layer}
  % Three service modules extracted to decouple the API routes from
  % infrastructure concerns:
  %
  % pipeline_service.py (run_pipeline_sync): downloads the SAS file via
  %   blob_service, executes the 8-node LangGraph pipeline, cleans up
  %   the temp file. Called by queue_service consumer or BackgroundTasks.
  %
  % blob_service.py (BlobStorageService): upload / download_to_temp /
  %   list_files / delete_folder. Targets Azure Blob Storage;
  %   falls back to local disk (backend/uploads/) when
  %   AZURE_STORAGE_CONNECTION_STRING is absent.
  %
  % queue_service.py (PipelineQueueService): enqueues JSON jobs into
  %   Azure Queue Storage (base64-encoded, visibility_timeout=300 s).
  %   Daemon consumer thread polls every 2 s and calls run_pipeline_sync.
  %   Poison-message guard: dead-letters after 5 dequeue attempts.
  %   Falls back to BackgroundTasks when not configured.

  \subsection{Containerisation and CI/CD}
  \label{subsec:docker}
  % Docker Compose: redis:7-alpine + backend + frontend (nginx).
  % GitHub Actions: 6-job pipeline — lint, test, benchmark,
  %   docker-build, CodeQL, azure-deploy.
  % Azure deploy: az containerapp update --image (image-only update;
  %   secrets untouched). Replaces az containerapp up which reset
  %   all secrets on every deploy.

  \subsection{Azure Key Vault Integration}
  \label{subsec:keyvault}
  % azure_setup.sh provisions secrets in Key Vault and grants the
  % Container App's managed identity the Key Vault Secrets User role.
  % Container App references secrets via --secrets keyvaultref: URIs.
  % Secrets added: OLLAMA_API_KEY, OLLAMA_BASE_URL, FRONTEND_URL
  % (CORS), GROQ_API_KEY, CODARA_JWT_SECRET, AZURE_OPENAI_API_KEY.
  % CORS_ORIGINS and OLLAMA_MODEL env vars added to Container App.

\section{Frontend Dashboard}
\label{sec:frontend-deployment}

  \subsection{Technology Stack and Architecture}
  \label{subsec:frontend-stack}

  \subsection{Key Interface Views}
  \label{subsec:workspace-ui}

  \subsection{End-to-End Integration Flow}
  \label{subsec:e2e-integration}

\section{Conclusion}
\label{sec:ch6-conclusion}


% =============================================================================
\chapter{Evaluation, Results, and Discussion}
\label{ch:evaluation}

\section{Introduction}
\label{sec:ch7-intro}

\section{Evaluation Methodology}
\label{sec:eval-methodology}

  \subsection{Evaluation Framework and Metrics}
  \label{subsec:metrics}

  \subsection{Experimental Protocol}
  \label{subsec:experimental-protocol}

\section{Boundary Detection Results}
\label{sec:boundary-results}

  \subsection{Deterministic Detector Performance}
  \label{subsec:det-perf}

  \subsection{LLM Fallback Contribution}
  \label{subsec:llm-boundary-contribution}

  \subsection{Discussion: Gap to Target}
  \label{subsec:boundary-gap}

\section{RAPTOR Ablation Study}
\label{sec:raptor-results}

  \subsection{Experimental Setup}
  \label{subsec:ablation-setup}

  \subsection{Retrieval Quality}
  \label{subsec:retrieval-quality}

  \subsection{Translation Quality Impact}
  \label{subsec:raptor-translation-impact}

  \subsection{HyperRAPTOR Pilot Results}
  \label{subsec:hyperraptor-results}

\section{Complexity Scoring Calibration}
\label{sec:complexity-results}

  \subsection{Calibration Curve and ECE}
  \label{subsec:ece-results}

  \subsection{Risk-Level Distribution and Routing Statistics}
  \label{subsec:risk-distribution}

\section{End-to-End Translation Results}
\label{sec:translation-results}

  \subsection{Overall Translation Success Rate}
  \label{subsec:overall-success}

  \subsection{Torture Test: 10/10 Success}
  \label{subsec:torture-test-results}
  % 10-block torture test (RETAIN + BY-group FIRST./LAST., missing value
  % logic, PROC SQL correlated subquery, Macro + \%DO loop, PROC MEANS
  % CLASS OUTPUT, PROC SORT NODUPKEY, hash object, multi-level nested
  % macro, PROC TRANSPOSE, complex WHERE + FORMAT).
  % Result: 10/10 SUCCESS, 3 Z3 formal proofs, 0 counterexamples.
  % Total time: 104.2 s. Model: Ollama nemotron-3-super:cloud (Azure
  % circuit-breaker open locally during test run).

  \subsection{Z3 Verification Outcomes}
  \label{subsec:z3-outcomes}

  \subsection{SemantiCheck Results}
  \label{subsec:semanticheck-results}
  % SemantiCheck (L3+L4 only; L1/L2 standalone eval) on 10-block
  % torture test:
  %
  % Avg SCS: 0.552 | VERIFIED or LIKELY\_CORRECT: 5/10
  % Avg L3 (STG): 0.508 | Avg L4 (Oracle): 0.595
  %
  % Notable findings:
  % - RETAIN + BY-group: SCS 0.900 (VERIFIED) — L3=1.00, L4=0.80
  % - PROC SORT NODUPKEY: SCS 0.925 (VERIFIED) — L3=1.00, L4=0.85
  % - PROC TRANSPOSE: SCS 0.150 (LIKELY\_INCORRECT) — translation
  %   executes but does not preserve structural intent.
  % - Multi-level macro: SCS 0.275 (LIKELY\_INCORRECT) — macros remain
  %   the primary semantic failure class.
  % Key finding: CodeBLEU would not have flagged TRANSPOSE or macro
  % translations as wrong; SemantiCheck catches structural intent loss
  % that execution-only testing misses.

  \subsection{Failure Mode Analysis}
  \label{subsec:failure-mode-analysis}

  \subsection{Streaming Performance}
  \label{subsec:streaming-perf}

\section{System Performance and LLM Cost}
\label{sec:system-perf}

  \subsection{API Latency}
  \label{subsec:api-latency}

  \subsection{LLM Token Cost}
  \label{subsec:llm-cost}

  \subsection{Test Suite Coverage}
  \label{subsec:test-coverage}

\section{Strengths of the Proposed Framework}
\label{sec:strengths}

  \subsection{Translation Quality and Retrieval Gains}
  \label{subsec:strength-success}

  \subsection{Z3 Formal Verification Catches Real Semantic Errors}
  \label{subsec:strength-z3}

  \subsection{SemantiCheck: Beyond Execution Correctness}
  \label{subsec:strength-semanticheck}
  % SemantiCheck identifies structural intent loss that passes both
  % syntax and execution checks. The composite SCS provides a
  % calibrated confidence signal independent of the generating LLM.

  \subsection{CDAIS: Formal Coverage Guarantees for Adversarial Patterns}
  \label{subsec:strength-cdais}

  \subsection{Modularity and Production-Ready Deployment}
  \label{subsec:strength-modularity}

\section{Limitations and Critical Assessment}
\label{sec:limitations}

  \subsection{Boundary Accuracy Below Target (79.3\% vs.\ 90\%)}
  \label{subsec:limit-boundary}

  \subsection{Macro Meta-Programming: Primary Failure Mode}
  \label{subsec:limit-macros}
  % SemantiCheck confirms macros as the primary semantic failure class:
  % LIKELY\_INCORRECT on both macro-heavy torture blocks (SCS 0.275).
  % L3 score = 0.00 for Macro + \%DO (no STG operations matched),
  % indicating the translation structurally diverges from the intent.

  \subsection{SemantiCheck L3/L4 Calibration Limitations}
  \label{subsec:limit-sc-calibration}
  % L4 oracle relies on LLM-simulated SAS execution: the LLM may
  % mispredict complex SAS semantics (e.g.\ PROC TRANSPOSE reshaping).
  % L3 STG extraction uses regex heuristics on SAS source — does not
  % parse the full SAS grammar. Both layers have known blind spots for
  % deeply nested macros and IML procedures.

  \subsection{Operational and Scope Limitations}
  \label{subsec:limit-operational}

\section{Integration with Design Science Research}
\label{sec:dsr-integration}

\section{Comparison with Related Work}
\label{sec:comparison}

  \subsection{vs.\ Industrial Migration Tools}
  \label{subsec:vs-industrial}

  \subsection{vs.\ Academic Code Translation Approaches}
  \label{subsec:vs-academic}

  \subsection{Novelty Claims}
  \label{subsec:novelty}
  % Five distinct novelty claims:
  % 1. 3-tier RAG with RAPTOR hierarchical clustering for SAS migration.
  % 2. Z3 SMT formal verification of generated Python code (4 patterns).
  % 3. CDAIS: Z3-driven minimum-witness adversarial test synthesis with
  %    formal coverage certificates for 6 SAS semantic error classes.
  % 4. SemantiCheck: 4-layer composite semantic verification (Formal +
  %    Behavioural + Contract/STG + Oracle/LLM) producing a calibrated
  %    SCS independent of the generating LLM.
  % 5. LLM-as-SAS-oracle (L4): behavioural testing without a SAS licence,
  %    using the LLM to simulate SAS execution on synthesised inputs.

\section{Future Work and Research Directions}
\label{sec:future-work}

  \subsection{Short-Term Improvements}
  \label{subsec:future-short-term}
  % - Integrate SemantiCheck L1/L2 into the main pipeline
  %   (currently standalone eval script only).
  % - Implement CDAIS synthesizer.py and coverage_oracle.py modules
  %   (architecture designed; constraint_catalog.py complete).
  % - Expand KB: 330 → 380 pairs (target: weak macro/TRANSPOSE categories).
  % - Frontend deployment (Azure Static Web Apps or second Container App).

  \subsection{Medium-Term Research Directions}
  \label{subsec:future-mid-term}
  % - Full SAS grammar parser for L3 STG extraction (replace heuristics).
  % - Fine-tuned local model (QLoRA on Qwen2.5-Coder) for Tier 0 offline.
  % - HyperRAPTOR (Poincaré ball) full evaluation on MOD/HIGH partitions.
  % - CDAIS extension to JOIN type and LAG queue failure classes.

  \subsection{Long-Term Scalability and Integration}
  \label{subsec:future-long-term}

\section{Conclusion}
\label{sec:ch7-conclusion}


% =============================================================================
\chapter*{General Conclusion}
\addcontentsline{toc}{chapter}{General Conclusion}
\markboth{General Conclusion}{General Conclusion}
\thispagestyle{fancy}

% =============================================================================
\bibliography{references}

\chapter*{Webography}
\addcontentsline{toc}{chapter}{Webography}
\thispagestyle{plain}
\begin{itemize}
  \item FastAPI documentation. \url{https://fastapi.tiangolo.com/} [Accessed: April 2026]
  \item LangGraph documentation. \url{https://langchain-ai.github.io/langgraph/} [Accessed: April 2026]
  \item LanceDB documentation. \url{https://lancedb.github.io/lancedb/} [Accessed: April 2026]
  \item Z3 Theorem Prover. \url{https://github.com/Z3Prover/z3} [Accessed: April 2026]
  \item Nomic Embed v1.5. \url{https://www.nomic.ai/blog/posts/nomic-embed-text-v1} [Accessed: April 2026]
  \item SAS Language Reference. \url{https://support.sas.com/documentation/} [Accessed: April 2026]
  \item Azure Blob Storage. \url{https://learn.microsoft.com/azure/storage/blobs/} [Accessed: April 2026]
  \item Azure Queue Storage. \url{https://learn.microsoft.com/azure/storage/queues/} [Accessed: April 2026]
  \item Weiser, M. (1984). Program slicing. \textit{IEEE Transactions on Software Engineering}, 10(4), 352--357.
\end{itemize}

\appendix
\renewcommand{\thechapter}{\Alph{chapter}}
\titleformat{\chapter}[display]
  {\normalfont\Large\bfseries\centering}
  {Appendix \thechapter}{10pt}{\LARGE}

\chapter{Gold Standard Corpus Sample}
\label{app:gold-corpus}

\chapter{RAPTOR Tree Visualisation}
\label{app:raptor-tree}

\chapter{Z3 SMT Pattern Listings}
\label{app:z3-patterns}

\chapter{Torture Test --- Annotated SAS Fixture}
\label{app:torture-test}

\chapter{SemantiCheck Score Breakdown --- Torture Test}
\label{app:semanticheck-scores}

\chapter{CDAIS Constraint Catalogue}
\label{app:cdais-catalogue}

\chapter{API Endpoint Reference}
\label{app:api-reference}

% =============================================================================
\end{document}
